\section{Sistemas y aplicaciones similares}
\label{sec:similares}

Antes de fijar el alcance del proyecto se relevaron las herramientas
disponibles en el mercado que abordan, total o parcialmente, el problema
de la preparación para entrevistas laborales. El relevamiento se centró
en cinco productos representativos del estado del arte, seleccionados
por cobertura de prensa y tracción reportada en el último año, y
distribuidos entre las dos categorías observables: simuladores
asistidos por IA (Interview Warmup, Big Interview, Yoodli, Final Round
AI) y plataformas de práctica entre pares (Pramp). Se descartaron del
relevamiento las plataformas de evaluación corporativa
(p.\,ej.\ HireVue) por estar orientadas a reclutadores y no al
candidato.

\subsection{Productos relevados}

\paragraph{Interview Warmup (Google).} Producto gratuito orientado a
candidatos de ingreso a empresas tecnológicas. Presenta preguntas en
texto y permite respuesta hablada por el navegador. Genera una
visualización de patrones lingüísticos (palabras frecuentes,
referencias profesionales, muletillas) sobre la transcripción. No
analiza video ni emite un puntaje cuantitativo y no contempla
adaptación dinámica de la dificultad ni interacción con otros usuarios.

\paragraph{Big Interview.} Suite comercial con planes individuales e
institucionales. Combina contenidos formativos extensos (videos, guías,
plantillas STAR) con un módulo de práctica donde el usuario se graba a
sí mismo respondiendo preguntas y recibe un panel posterior con
métricas básicas (\textit{eye contact}, \textit{filler words}, ritmo).
La retroalimentación se entrega como \textit{dashboard} lateral
clásico y no en tiempo real durante la entrevista. La adaptación al
candidato es limitada y depende de la curaduría manual de bancos de
preguntas por industria.

\paragraph{Yoodli.} Aplicación enfocada en la fluidez del discurso
hablado más que en el contenido técnico de la entrevista. Soporta
sesiones grabadas y modo en vivo durante reuniones de
videoconferencia, con métricas prosódicas (palabras por minuto, pausas,
muletillas, riqueza léxica) entregadas al cierre. Su análisis no
verbal se limita a la voz y no incorpora análisis postural ni de
contacto visual; tampoco habilita sesiones colaborativas entre pares.

\paragraph{Final Round AI.} Producto reciente posicionado como
\textit{interview copilot}: durante una entrevista real con un
reclutador humano, sugiere al candidato respuestas o estructuras
argumentales en una ventana lateral. El modo \textit{mock} es
secundario y replica el patrón pregunta-respuesta con feedback
posterior. El énfasis del producto está en la asistencia en vivo y no
en el entrenamiento longitudinal con plan de mejora; el componente
ético del uso del \textit{copilot} en entrevistas reales está aún en
discusión pública.

\paragraph{Pramp.} Plataforma colaborativa donde dos candidatos se
emparejan por compatibilidad de rol y se entrevistan mutuamente, sin
intervención de IA. Cada participante alterna roles candidato y
entrevistador, y se evalúan con una rúbrica predefinida. Su valor
reside en la práctica humana realista, pero la calidad del feedback
queda íntegramente en manos del par y carece de métricas objetivas
multimodales.

\subsection{Comparación funcional}

La tabla \ref{tab:similares} resume el cruce de capacidades observadas
con los ejes que estructuran este proyecto: análisis multimodal,
\textit{feedback} en tiempo real, adaptación al candidato, modalidad
colaborativa entre pares, gamificación y atención explícita a la
accesibilidad.

\begin{table}[ht]
\centering
\renewcommand{\arraystretch}{1.25}
\footnotesize
\begin{tabularx}{\linewidth}{@{}l Y Y Y Y Y c@{}}
\toprule
\textbf{Producto} &
\textbf{Análisis voz} &
\textbf{Análisis video} &
\textbf{Feedback en vivo} &
\textbf{Adaptación al nivel} &
\textbf{Práctica entre pares} &
\textbf{WCAG AA} \\
\midrule
Interview Warmup & Limitado (texto) & No & No & Limitada & No & Parcial \\
Big Interview    & Sí (post)        & Sí (post) & No & Limitada & No & Parcial \\
Yoodli           & Sí (vivo+post)   & No & Parcial & No & No & Parcial \\
Final Round AI   & Sí (vivo)        & No & Sí (sugerencias) & Limitada & No & No documentado \\
Pramp            & No               & No (peer humano) & N/A & N/A & Sí (núcleo) & No documentado \\
\midrule
\textbf{Este proyecto} &
\textbf{Sí (vivo)} &
\textbf{Sí (vivo)} &
\textbf{Sí (aura)} &
\textbf{Sí (IRT)} &
\textbf{Sí (peer mock)} &
\textbf{Sí (AA explícito)} \\
\bottomrule
\end{tabularx}
\caption{Comparación funcional con cinco productos representativos.
La fila final corresponde al simulador propuesto.}
\label{tab:similares}
\end{table}

\subsection{Posicionamiento y diferenciación}

El relevamiento permite ubicar el proyecto en una intersección poco
cubierta del espacio de soluciones existentes. La oferta actual se
distribuye en tres polos: (i) simuladores con \textit{feedback}
posterior y análisis incompleto del lenguaje no verbal (Big Interview,
Yoodli), (ii) asistentes en vivo orientados a entrevistas reales con
componente ético cuestionado (Final Round AI), y (iii) plataformas de
práctica entre pares sin métricas objetivas (Pramp). Ninguna de las
herramientas relevadas combina simultáneamente: análisis multimodal
sincrónico de voz y video con cómputo en cliente, retroalimentación
visual integrada al avatar (no como \textit{dashboard} lateral),
adaptación de dificultad por modelo de habilidad por competencia,
modalidad colaborativa entre pares con observador instrumentado, y
conformidad explícita con WCAG 2.2 nivel AA con rutas alternativas
para usuarios con discapacidad.

La propuesta del presente proyecto se diferencia en cuatro ejes
verificables. Primero, el \textit{aura luminosa procedural} alrededor
del entrevistador-avatar reemplaza al panel de métricas tradicional y
materializa la metáfora del \textit{dojo}, sosteniendo la inmersión
durante la sesión (cubre \rf{05} y \rf{15}). Segundo, la totalidad del
análisis de video ocurre en el navegador mediante MediaPipe, sin envío
de imagen cruda a servicios externos, lo que mejora simultáneamente la
privacidad y el costo operativo (\rnf{05}). Tercero, la sesión
\textit{peer mock} sobre WebRTC habilita observación instrumentada
con métricas cruzadas y comentarios anclados a \textit{timestamps},
una capacidad ausente en Pramp y en los productos asistidos por IA
relevados (\rf{10}--\rf{12}). Cuarto, la accesibilidad se trata como
requerimiento de primer nivel y no como adición tardía: todos los
canales del aura tienen sustituto auditivo o textual y la operación
total del simulador es posible con teclado, conmutadores o voz
(\rf{17}--\rf{20} y \rnf{08}).

Este posicionamiento orienta también el alcance del MVP: el equipo
prioriza completar las capacidades diferenciadoras antes de replicar
funciones ya bien resueltas por la competencia (p.\,ej.\ catálogos
extensos de preguntas curados manualmente, formación en redacción de
\textit{currícula}, integración con plataformas de empleabilidad
externas), las cuales quedan documentadas como extensiones futuras.
